\section{Introduction}\label{sec:intro}

Bidding into wholesale electricity markets is a natural task to delegate to software. Day-ahead and intraday auctions clear every hour or every quarter hour, the rules are public, and the feedback a bidder needs, the clearing price and its own dispatched volume, arrives promptly. These are the conditions under which reinforcement learning is easy to deploy, and under which competition authorities worry about it. The concern, raised by the \citet{oecd2017} and developed by \citet{harrington2018}, is that independent learning algorithms may arrive at supra-competitive prices without any agreement, leaving competition law with no obvious handle on the result. \citet{calvano2020} gave the concern an experimental basis: independent Q-learners in a repeated Bertrand game learn prices well above the static Nash level and, when one deviates, punish and then forgive, which is the fingerprint of a reward-punishment scheme. \citet{seredynski2026} report cases of the same pattern for learning agents bidding into a networked electricity market.

The difficulty is that high prices are not collusion. Tacit collusion in the sense of the folk theorem \citep{fudenberg1986} means that each firm holds back because it expects retaliation if it does not. Many other things can also produce high prices: a coarse action grid, a market rule under which a lone undercut does not move the price, an exploration schedule that dies before the competitive race has finished, or a value estimate that was never revisited. \citet{abada2023} and \citet{abada2024} show that seemingly collusive outcomes in learning simulations can come from the learners' own limitations rather than from strategy, and \citet{banchio2022} show that the auction rule itself changes what independent learners settle on. A regulator who sees only prices cannot tell these cases apart. The question this paper asks is whether the indicators used in the simulation literature can.

The setting is deliberately small. Three symmetric firms with 60~MW of capacity each and a marginal cost of \eur{10}/MWh bid one price from a 15-rung grid between \eur{10} and \eur{100}/MWh into a uniform-price auction with inelastic demand of 100~MW. Each firm is an independent tabular Q-learner that observes only last round's bids and its own profit, with no communication and no shared parameters. Because any two firms cover demand, the clearing price is always the second-lowest bid, so from a profile at which the two lowest bids are tied, which includes every symmetric profile, a lone undercutter is paid the rivals' price for its full 60~MW and the price does not move.

First, the learners reach supra-competitive prices. At the pre-registered horizon of 1M rounds the collusion index $\Dl$ averages 0.375 (95\% CI 0.354--0.397) and no seed has converged; at 5M rounds every seed has converged and $\Dl$ is 0.621 (0.612--0.630). Second, a forced one-round undercut is answered but rarely followed by a return: with firm 1 as deviator, rivals' bids fall in 91\% of converged seeds and are back at their old level within 24 rounds in only 21\%, and the per-seed verdict changes with the deviator's identity in 61 of 100 seeds. Third, the shortsightedness ablation fails. At 5M rounds a learner with no future ($\gamma = 0$) settles at $\Dl = 0.370$, removing about 40\% of the full learner's price premium but leaving $\Dl$ well above the pre-registered collapse threshold of 0.25, and a learner with no memory of the past reaches $\Dl = 0.943$, higher than the full learner. Fourth, the unclaimed-temptation fingerprint, a profitable deviation left unplayed, is present in 100\% of seeds, but it is present in 100\% of memory-0 control seeds too, and a memory-0 learner cannot anticipate anything. By the indicators' own rules, the evidence does not support attributing the high prices to tacit collusion. This is a negative result and it is the main finding.

The memory-0 price is then the puzzle, and a second set of five pre-registered experiments addresses it. Two of the five predictions failed, and the failures are informative. The memory-0 price was predicted to fall when exploration decays ten times more slowly; it does not. It was predicted to fall a lot under pay-as-bid; it falls by 0.10. It also returns after a perturbation that restarts exploration from $\varepsilon = 0.5$ and swaps the low-bidder role in 70\% of seeds. What the experiments support instead is a mechanism, which we propose as a verbal argument rather than a theorem: an Edgeworth-type cycle under $\varepsilon$-greedy exploration, in which the price walks down the grid one rung per learning event and jumps back to the top in a single event, so the cycle should live near the cap, and the less noise there is, the more time it should spend there. A constant-$\varepsilon$ occupation measure is consistent with the frozen price being the small-noise limit of the exploring process for both learners: the memory-0 time average is within 0.01 of its frozen value at $\varepsilon = 0.001$, while the memory-1 average meets its frozen value only at $\varepsilon = 10^{-4}$, roughly 30 times less noise, and only one point was run there. Finally, a Bellman-residual test on the Q-tables separates the two learners. Both learners' deviation estimates are stale, but the share of temptations where one fresh backup would already flip the decision is exactly 100\% in every memory-0 seed and at most 80\% in any full-learner seed, with a mean of 31\%. In the other 69\% of the full learner's temptations the table's own value of the state reached by deviating is low enough that one fresh backup would leave the deviation unattractive; that backup takes its continuation from the current table, which may itself be stale. That is a learned continuation, and the two distributions do not overlap.

The contribution is a replication-plus-ablation study with a negative result, a mechanism, and a diagnostic. The negative result is that the behavioural indicators of tacit collusion (punishment, ablation and unclaimed temptation) cannot distinguish a learner that has learned a continuation from one that has stopped looking, in a market where a lone deviation does not move the price: the ablation and the temptation test read the same for both learners, and the punishment test separates them only because a one-state table has nothing to react to. The mechanism is the Edgeworth-type cycle and its small-noise limit; the argument is verbal and supported by experiments, not by a theorem. The diagnostic is the Bellman-residual share, which needs the Q-table and is therefore a simulation-side tool, not a screen for real bid data. The paper offers no theory of algorithmic collusion and never claims that the agents collude.

Section~\ref{sec:related} places the paper in the literature on algorithmic pricing and on electricity auctions. Section~\ref{sec:model} sets out the market, the learners and the four tests, and Section~\ref{sec:design} the experimental design and acceptance rules. Section~\ref{sec:results} reports price levels, punishment, the ablation and the temptation test. Section~\ref{sec:phase1} reports the five follow-up experiments and the Bellman-residual test. Section~\ref{sec:discussion} interprets the results and lists their limits; Section~\ref{sec:conclusion} concludes. Reproduction details are in Appendix~\ref{app:repro}.
